\section{Holm Bonferroni - For all components touched for all commits} 
\label{sec:codebase_graph_latex.repository_graph.component_time_series.holm_bonferroni} 
A Pairwise Post-hoc Analysis with Holm-Bonferroni Correction table for all components touched on average. 

\paragraph{Post-hoc Pairwise Analysis (Holm-Bonferroni):}
Following the multi-factor significance analysis, non-parametric pairwise comparisons 
were executed to isolate directional stochastic dominance patterns between groups. 
To counteract family-wise Type I error inflation across multiple group permutations 
without experiencing the severe statistical power depletion of a traditional 
Bonferroni adjustment, a sequentially rejective Holm-Bonferroni step-down 
correction is applied.
\begin{table}[H]
\centering
\setlength{\tabcolsep}{2pt}
\caption{Significant Holm-Bonferroni Pairwise results for packages and number of commits}
\resizebox{\textwidth}{!}{% 
\begin{tabular}{l r r l l r r}
\toprule
\textbf{Component} & \textbf{Commits} & \makecell[t]{\textbf{Number of}\\\textbf{Developers}} & \multicolumn{2}{c}{\textbf{Group}} & \textbf{Mean Diff} & \makecell[t]{\textbf{Adj $P$}\\\textbf{Value}} \\ 
\cmidrule(lr){4-5}
 &  &  & \textbf{A} & \textbf{B} &  &  \\ 
\midrule
packages & 1 & 16434 & Founder Moderate & Founder Transient & 1.39 & 0.0000 \\ 
packages & 1 & 16434 & Founder Moderate & Later Joiner Transient & 1.44 & 0.0000 \\ 
packages & 1 & 16434 & Founder Sustained & Founder Transient & 1.01 & 0.0000 \\ 
packages & 1 & 16434 & Founder Sustained & Later Joiner Transient & 1.06 & 0.0000 \\ 
packages & 1 & 16434 & Founder Transient & Later Joiner Moderate & -1.40 & 0.0000 \\ 
packages & 1 & 16434 & Founder Transient & Later Joiner Sustained & -1.21 & 0.0000 \\ 
packages & 1 & 16434 & Later Joiner Moderate & Later Joiner Transient & 1.45 & 0.0000 \\ 
packages & 1 & 16434 & Later Joiner Sustained & Later Joiner Transient & 1.26 & 0.0000 \\ 
packages & 2 & 3085 & Later Joiner Moderate & Later Joiner Transient & 0.53 & 0.0000 \\ 
packages & 5 & 1362 & Later Joiner Moderate & Later Joiner Transient & 0.53 & 0.0214 \\ 
\bottomrule
\end{tabular}%
} 
\label{tab:codebase_graph_latex_repository_graph_component_time_series_holm_bonferroni_sig} 
\end{table}
\newpage 
\begin{table}[H]
\centering
\setlength{\tabcolsep}{2pt}
\caption{Significant Holm-Bonferroni Pairwise results for classes and number of commits}
\resizebox{\textwidth}{!}{% 
\begin{tabular}{l r r l l r r}
\toprule
\textbf{Component} & \textbf{Commits} & \makecell[t]{\textbf{Number of}\\\textbf{Developers}} & \multicolumn{2}{c}{\textbf{Group}} & \textbf{Mean Diff} & \makecell[t]{\textbf{Adj $P$}\\\textbf{Value}} \\ 
\cmidrule(lr){4-5}
 &  &  & \textbf{A} & \textbf{B} &  &  \\ 
\midrule
classes & 1 & 16434 & Founder Moderate & Founder Transient & 6.26 & 0.0001 \\ 
classes & 1 & 16434 & Founder Moderate & Later Joiner Transient & 6.73 & 0.0000 \\ 
classes & 1 & 16434 & Founder Sustained & Founder Transient & 1.82 & 0.0001 \\ 
classes & 1 & 16434 & Founder Sustained & Later Joiner Transient & 2.30 & 0.0000 \\ 
classes & 1 & 16434 & Founder Transient & Later Joiner Moderate & -1.33 & 0.0000 \\ 
classes & 1 & 16434 & Founder Transient & Later Joiner Sustained & -3.87 & 0.0006 \\ 
classes & 1 & 16434 & Founder Transient & Later Joiner Transient & 0.47 & 0.0001 \\ 
classes & 1 & 16434 & Later Joiner Moderate & Later Joiner Transient & 1.80 & 0.0000 \\ 
classes & 1 & 16434 & Later Joiner Sustained & Later Joiner Transient & 4.34 & 0.0000 \\ 
classes & 2 & 3085 & Founder Transient & Later Joiner Sustained & -0.75 & 0.0198 \\ 
classes & 2 & 3085 & Founder Transient & Later Joiner Transient & 2.70 & 0.0002 \\ 
classes & 2 & 3085 & Later Joiner Moderate & Later Joiner Transient & 0.91 & 0.0001 \\ 
classes & 5 & 1362 & Later Joiner Moderate & Later Joiner Transient & 1.01 & 0.0340 \\ 
\bottomrule
\end{tabular}%
} 
\label{tab:codebase_graph_latex_repository_graph_component_time_series_holm_bonferroni_sig} 
\end{table}
\newpage 
\begin{table}[H]
\centering
\setlength{\tabcolsep}{2pt}
\caption{Significant Holm-Bonferroni Pairwise results for methods and number of commits}
\resizebox{\textwidth}{!}{% 
\begin{tabular}{l r r l l r r}
\toprule
\textbf{Component} & \textbf{Commits} & \makecell[t]{\textbf{Number of}\\\textbf{Developers}} & \multicolumn{2}{c}{\textbf{Group}} & \textbf{Mean Diff} & \makecell[t]{\textbf{Adj $P$}\\\textbf{Value}} \\ 
\cmidrule(lr){4-5}
 &  &  & \textbf{A} & \textbf{B} &  &  \\ 
\midrule
methods & 1 & 16434 & Founder Moderate & Founder Transient & 21.43 & 0.0001 \\ 
methods & 1 & 16434 & Founder Moderate & Later Joiner Transient & 22.81 & 0.0000 \\ 
methods & 1 & 16434 & Founder Sustained & Founder Transient & 7.75 & 0.0003 \\ 
methods & 1 & 16434 & Founder Sustained & Later Joiner Transient & 9.13 & 0.0000 \\ 
methods & 1 & 16434 & Founder Transient & Later Joiner Moderate & -5.09 & 0.0000 \\ 
methods & 1 & 16434 & Founder Transient & Later Joiner Sustained & -21.24 & 0.0003 \\ 
methods & 1 & 16434 & Founder Transient & Later Joiner Transient & 1.38 & 0.0000 \\ 
methods & 1 & 16434 & Later Joiner Moderate & Later Joiner Transient & 6.47 & 0.0000 \\ 
methods & 1 & 16434 & Later Joiner Sustained & Later Joiner Transient & 22.62 & 0.0000 \\ 
methods & 2 & 3085 & Founder Transient & Later Joiner Sustained & -11.96 & 0.0052 \\ 
methods & 2 & 3085 & Founder Transient & Later Joiner Transient & 7.84 & 0.0000 \\ 
methods & 2 & 3085 & Later Joiner Moderate & Later Joiner Transient & 3.64 & 0.0000 \\ 
methods & 5 & 1362 & Later Joiner Moderate & Later Joiner Transient & 4.43 & 0.0086 \\ 
\bottomrule
\end{tabular}%
} 
\label{tab:codebase_graph_latex_repository_graph_component_time_series_holm_bonferroni_sig} 
\end{table}
\newpage 

\begin{table}[H]
\centering
\setlength{\tabcolsep}{2pt}
\caption{Not Significant Holm-Bonferroni Pairwise results for packages and number of commits}
\resizebox{\textwidth}{!}{% 
\begin{tabular}{l r r l l r r}
\toprule
\textbf{Component} & \textbf{Commits} & \makecell[t]{\textbf{Number of}\\\textbf{Developers}} & \multicolumn{2}{c}{\textbf{Group}} & \textbf{Mean Diff} & \makecell[t]{\textbf{Adj $P$}\\\textbf{Value}} \\ 
\cmidrule(lr){4-5}
 &  &  & \textbf{A} & \textbf{B} &  &  \\ 
\midrule
packages & 1 & 16434 & Founder Moderate & Founder Sustained & 0.38 & 1.0000 \\ 
packages & 1 & 16434 & Founder Moderate & Later Joiner Moderate & -0.01 & 1.0000 \\ 
packages & 1 & 16434 & Founder Moderate & Later Joiner Sustained & 0.18 & 1.0000 \\ 
packages & 1 & 16434 & Founder Sustained & Later Joiner Moderate & -0.39 & 1.0000 \\ 
packages & 1 & 16434 & Founder Sustained & Later Joiner Sustained & -0.20 & 1.0000 \\ 
packages & 1 & 16434 & Founder Transient & Later Joiner Transient & 0.05 & 1.0000 \\ 
packages & 1 & 16434 & Later Joiner Moderate & Later Joiner Sustained & 0.19 & 1.0000 \\ 
packages & 5 & 1362 & Founder Moderate & Later Joiner Moderate & -0.54 & 1.0000 \\ 
packages & 5 & 1362 & Founder Sustained & Later Joiner Sustained & -0.46 & 1.0000 \\ 
packages & 5 & 1362 & Founder Transient & Later Joiner Transient & -1.76 & 1.0000 \\ 
packages & 10 & 924 & Founder Moderate & Founder Sustained & 1.41 & 0.8267 \\ 
packages & 10 & 924 & Founder Moderate & Later Joiner Moderate & -1.61 & 0.1575 \\ 
packages & 10 & 924 & Founder Moderate & Later Joiner Sustained & -0.63 & 0.8267 \\ 
packages & 10 & 924 & Founder Sustained & Later Joiner Sustained & -2.03 & 0.3229 \\ 
\bottomrule
\end{tabular}%
} 
\end{table}
\newpage 
\begin{table}[H]
\centering
\setlength{\tabcolsep}{2pt}
\caption{Not Significant Holm-Bonferroni Pairwise results for classes and number of commits}
\resizebox{\textwidth}{!}{% 
\begin{tabular}{l r r l l r r}
\toprule
\textbf{Component} & \textbf{Commits} & \makecell[t]{\textbf{Number of}\\\textbf{Developers}} & \multicolumn{2}{c}{\textbf{Group}} & \textbf{Mean Diff} & \makecell[t]{\textbf{Adj $P$}\\\textbf{Value}} \\ 
\cmidrule(lr){4-5}
 &  &  & \textbf{A} & \textbf{B} &  &  \\ 
\midrule
classes & 1 & 16434 & Founder Moderate & Founder Sustained & 4.44 & 1.0000 \\ 
classes & 1 & 16434 & Founder Moderate & Later Joiner Moderate & 4.93 & 1.0000 \\ 
classes & 1 & 16434 & Founder Moderate & Later Joiner Sustained & 2.39 & 0.7561 \\ 
classes & 1 & 16434 & Founder Sustained & Later Joiner Moderate & 0.49 & 1.0000 \\ 
classes & 1 & 16434 & Founder Sustained & Later Joiner Sustained & -2.05 & 0.7561 \\ 
classes & 1 & 16434 & Later Joiner Moderate & Later Joiner Sustained & -2.54 & 0.7561 \\ 
classes & 5 & 1362 & Founder Transient & Later Joiner Transient & -6.89 & 0.9533 \\ 
classes & 10 & 924 & Founder Moderate & Founder Sustained & -7.66 & 1.0000 \\ 
classes & 10 & 924 & Founder Moderate & Later Joiner Moderate & 18.61 & 0.5792 \\ 
classes & 10 & 924 & Founder Moderate & Later Joiner Sustained & 12.95 & 0.5163 \\ 
classes & 10 & 924 & Founder Sustained & Later Joiner Moderate & 26.27 & 0.9912 \\ 
classes & 10 & 924 & Founder Sustained & Later Joiner Sustained & 20.61 & 0.9912 \\ 
classes & 10 & 924 & Later Joiner Moderate & Later Joiner Sustained & -5.66 & 1.0000 \\ 
classes & 20 & 237 & Founder Sustained & Later Joiner Sustained & 25.59 & 0.2857 \\ 
\bottomrule
\end{tabular}%
} 
\end{table}
\newpage 
\begin{table}[H]
\centering
\setlength{\tabcolsep}{2pt}
\caption{Not Significant Holm-Bonferroni Pairwise results for methods and number of commits}
\resizebox{\textwidth}{!}{% 
\begin{tabular}{l r r l l r r}
\toprule
\textbf{Component} & \textbf{Commits} & \makecell[t]{\textbf{Number of}\\\textbf{Developers}} & \multicolumn{2}{c}{\textbf{Group}} & \textbf{Mean Diff} & \makecell[t]{\textbf{Adj $P$}\\\textbf{Value}} \\ 
\cmidrule(lr){4-5}
 &  &  & \textbf{A} & \textbf{B} &  &  \\ 
\midrule
methods & 1 & 16434 & Founder Moderate & Founder Sustained & 13.68 & 1.0000 \\ 
methods & 1 & 16434 & Founder Moderate & Later Joiner Moderate & 16.34 & 1.0000 \\ 
methods & 1 & 16434 & Founder Moderate & Later Joiner Sustained & 0.19 & 0.6748 \\ 
methods & 1 & 16434 & Founder Sustained & Later Joiner Moderate & 2.66 & 1.0000 \\ 
methods & 1 & 16434 & Founder Sustained & Later Joiner Sustained & -13.49 & 0.9405 \\ 
methods & 1 & 16434 & Later Joiner Moderate & Later Joiner Sustained & -16.15 & 0.5367 \\ 
methods & 5 & 1362 & Founder Transient & Later Joiner Transient & -24.55 & 0.8554 \\ 
methods & 10 & 924 & Founder Moderate & Founder Sustained & 13.47 & 1.0000 \\ 
methods & 10 & 924 & Founder Moderate & Later Joiner Moderate & 119.34 & 0.6249 \\ 
methods & 10 & 924 & Founder Moderate & Later Joiner Sustained & 86.42 & 0.6249 \\ 
methods & 10 & 924 & Founder Sustained & Later Joiner Moderate & 105.87 & 0.7351 \\ 
methods & 10 & 924 & Founder Sustained & Later Joiner Sustained & 72.95 & 0.7351 \\ 
methods & 10 & 924 & Later Joiner Moderate & Later Joiner Sustained & -32.92 & 1.0000 \\ 
methods & 20 & 237 & Founder Sustained & Later Joiner Sustained & 106.14 & 0.5131 \\ 
\bottomrule
\end{tabular}%
} 
\end{table}
\newpage 


\subsubsection*{Post-hoc Pairwise Analysis: Holm-Bonferroni Correction}
To evaluate specific distribution differences between isolated developer interaction 
groups without violating normality or homoscedastic constraints, we implement a 
two-sample rank-order framework adjusted by the sequentially rejective Holm method.

The raw $P$-values for a set of $m$ concurrent hypotheses are sorted in ascending 
order:
\begin{equation}
    P_{(1)} \le P_{2} \le \dots \le P_{(m)}
\end{equation}

Each index $i$ ($1 \le i \le m$) is evaluated sequentially against a dynamically 
scaled significance threshold. The condition for rejecting the local null hypothesis 
$H_{(i)}$ is defined as:
\begin{equation}
    P_{(i)} \le \frac{\alpha}{m - i + 1}
\end{equation}

Where:
\begin{itemize}
    \item $\alpha$ is the nominal family-wise error criteria cutoff (configured to $0.05$).
    \item $m$ is the total number of simultaneous pairwise comparisons evaluated.
    \item $i$ is the ordered ranking position of the active sub-test segment.
\end{itemize}
The sequence terminates at the first instance where $P_{(i)} > \frac{\alpha}{m - i + 1}$, 
retaining all remaining null structures.

\newpage 
